% from ega3-1-fr.tex
\oldpage[0\textsubscript{III}]{61}
\begin{proposition}[12.3.4]
\label{0.12.3.4-fr}
\emph{Soit $f:X\to Y$ un morphisme plat d'espaces annelés, et soient
$\mathcal F$, $\mathcal G$ deux $\mathcal O_Y$-Modules.}
\begin{enumerate}
\item[{\rm(i)}] \emph{Il existe un homomorphisme de bifoncteurs
cohomologiques
\[
\label{0.12.3.4.1-fr}
\tag{12.3.4.1}
f^*(\mathcal{E}xt_{\mathcal O_Y}^{\bullet}(\mathcal F,\mathcal G))
\xrightarrow{\sim}
\mathcal{E}xt_{\mathcal O_X}^{\bullet}(f^*(\mathcal F),f^*(\mathcal G))
\]
se réduisant en degré $0$ à l'homomorphisme canonique
(\hyperref[0.4.4.6-fr]{0\textsubscript{I}, 4.4.6}).}

\item[{\rm(ii)}] \emph{Il existe un morphisme canonique de suites spectrales
\[
\label{0.12.3.4.2-fr}
\tag{12.3.4.2}
E(\mathcal F,\mathcal G)\longrightarrow
E(f^*(\mathcal F),f^*(\mathcal G))
\]
qui, pour les termes $E_2$, se réduit aux homomorphismes
\[
\label{0.12.3.4.3-fr}
\tag{12.3.4.3}
H^p(Y,\mathcal{E}xt_{\mathcal O_Y}^q(\mathcal F,\mathcal G))
\longrightarrow
H^p(X,\mathcal{E}xt_{\mathcal O_X}^q(f^*(\mathcal F),f^*(\mathcal G)))
\]
déduits de (\hyperref[0.12.3.4.1-fr]{12.3.4.1}) et de
(\hyperref[0.12.1.3.1-fr]{12.1.3.1}).}
\end{enumerate}
\end{proposition}

\begin{proof}
\begin{enumerate}
\item[{\rm(i)}]
Comme $f^*$ est un foncteur exact dans la catégorie des
$\mathcal O_Y$-Modules (\hyperref[0.6.7.2-fr]{0\textsubscript{I}, 6.7.2}),
les foncteurs
\[
\mathcal G\leadsto
f^*(\mathcal{H}om_{\mathcal O_Y}(\mathcal F,\mathcal G))
\quad\hbox{et}\quad
\mathcal G\leadsto
\mathcal{H}om_{\mathcal O_X}(f^*(\mathcal F),f^*(\mathcal G))
\]
sont exacts à gauche; on déduit canoniquement de
(\hyperref[0.4.4.6-fr]{0\textsubscript{I}, 4.4.6}) un homomorphisme de leurs
foncteurs dérivés.

Pour calculer ces derniers, on prend une résolution injective
$\mathcal L^\bullet=(\mathcal L^i)$ de $\mathcal G$, et on a donc des
morphismes
\[
\mathcal H^p(f^*(\mathcal{H}om_{\mathcal O_Y}
(\mathcal F,\mathcal L^\bullet)))
\longrightarrow
\mathcal H^p(\mathcal{H}om_{\mathcal O_X}
(f^*(\mathcal F),f^*(\mathcal L^\bullet)))
\]
de cohomologies de complexes de faisceaux. D'ailleurs, en vertu de
l'exactitude de $f^*$, on a
\[
\mathcal H^p(f^*(\mathcal{H}om_{\mathcal O_Y}
(\mathcal F,\mathcal L^\bullet)))
=f^*(\mathcal H^p(\mathcal{H}om_{\mathcal O_Y}
(\mathcal F,\mathcal L^\bullet)))
=f^*(\mathcal{E}xt_{\mathcal O_Y}^p(\mathcal F,\mathcal G))
\]
par définition. D'autre part, l'exactitude de $f^*$ entraîne que
$f^*(\mathcal L^\bullet)$ est une résolution de $f^*(\mathcal G)$; si
$\mathcal L'^\bullet=(\mathcal L'^{i})$ est une résolution injective de
$f^*(\mathcal G)$ dans la catégorie des $\mathcal O_X$-Modules, il y a donc
un homomorphisme de complexes
$f^*(\mathcal L^\bullet)\to\mathcal L'^\bullet$, déterminé à une homotopie
près, et qui définit par suite un homomorphisme bien déterminé en cohomologie;
composant cet homomorphisme avec l'homomorphisme défini plus haut, on obtient
(\hyperref[0.12.3.4.1-fr]{12.3.4.1}).

\item[{\rm(ii)}]
Avec les notations précédentes, on a un homomorphisme de complexes de
faisceaux de $\mathcal O_X$-Modules
\[
f^*(\mathcal{H}om_{\mathcal O_Y}(\mathcal F,\mathcal L^\bullet))
\longrightarrow
\mathcal{H}om_{\mathcal O_X}(f^*(\mathcal F),\mathcal L'^\bullet).
\]
Soit $\mathcal M^{\bullet\bullet}$ une résolution injective de
Cartan-Eilenberg du complexe
$\mathcal{H}om_{\mathcal O_Y}(\mathcal F,\mathcal L^\bullet)$ dans la
catégorie des $\mathcal O_Y$-Modules; alors, en vertu de l'exactitude du
foncteur $f^*$, $f^*(\mathcal M^{\bullet\bullet})$ est une résolution de
Cartan-Eilenberg du complexe
$f^*(\mathcal{H}om_{\mathcal O_Y}(\mathcal F,\mathcal L^\bullet))$; si
$\mathcal M'^{\bullet\bullet}$ est une résolution injective de
Cartan-Eilenberg du complexe
$\mathcal{H}om_{\mathcal O_X}(f^*(\mathcal F),\mathcal L'^\bullet)$, il y a
donc (\hyperref[0.11.4.2-fr]{11.4.2}), un homomorphisme (déterminé à
homotopie près)
\[
f^*(\mathcal M^{\bullet\bullet})\longrightarrow
\mathcal M'^{\bullet\bullet}
\]
compatible avec l'homomorphisme considéré ci-dessus, autrement dit un
$f$-morphisme $\mathcal M^{\bullet\bullet}\to
\mathcal M'^{\bullet\bullet}$ de bicomplexes de faisceaux. On en déduit un
di-homomorphisme
$\Gamma(Y,\mathcal M^{\bullet\bullet})\to
\Gamma(X,\mathcal M'^{\bullet\bullet})$ de bicomplexes de modules, déterminé
à homotopie près, et un morphisme bien déterminé de suites spectrales
(\hyperref[0.11.3.2-fr]{11.3.2}), qui n'est autre que le morphisme
(\hyperref[0.12.3.4.2-fr]{12.3.4.2}) cherché, la caractérisation de
(\hyperref[0.12.3.4.3-fr]{12.3.4.3}) se déduisant aussitôt des définitions.
\end{enumerate}
\end{proof}

\begin{proposition}[12.3.5]
\label{0.12.3.5-fr}
\emph{Sous les hypothèses de (\hyperref[0.12.3.4-fr]{12.3.4}), supposons en
outre le faisceau d'anneaux $\mathcal O_Y$ cohérent; alors, pour tout
$\mathcal O_Y$-Module cohérent $\mathcal F$, les homomorphismes canoniques
(\hyperref[0.12.3.4.1-fr]{12.3.4.1}) sont bijectifs.}
\end{proposition}

\begin{proof}
